\section{Introduction}

RCOUNT packages are meant to be public evidence, not office-only artifacts.
Anyone who receives the package should be able to run the verifier and recover
the same transcript that an election office would see. For full-election
packages, that goal creates an ordinary engineering pressure: independent checks
over many public records should eventually run in parallel.

Parallel verification is only useful if it does not change what is being
verified. A faster verifier that reorders reports, hides a failing equation, or
returns a coarser pass/fail result would weaken the audit surface. The public
object is the transcript: the ordered list of equations that passed, plus the
same error behavior when a package fails.

This paper documents the first RCOUNT parallel verifier slice. The serial entry
point is \texttt{verify\_package}; the parallel entry point is
\texttt{verify\_package\_parallel}. The implementation uses Rayon through
\texttt{rayon::prelude} and parallelizes the contest selection-sum family first.
The research claim is deliberately narrow: the parallel verifier preserves the
serial transcript semantics for the implemented surface. It does not yet claim
measured speedups.
